\clearpage
\section{课件补充知识点（非已确认必考）}
本节补足课件中容易遗漏的知识点，配自编小例题。\textbf{出现在课件中，不代表老师已经确认它必考。}时间紧张时，先掌握正文的分辨率、灰度变换、均衡化与空间滤波，再读本节。

\subsection{人眼视觉：适应、辨别与相对亮度}
\textit{来源：Perception 课件第3--19页；教材\S2.1，印刷页48--54（PDF页50--56）。}
\begin{itemize}
\item \textbf{视锥细胞与视杆细胞：}视锥主要负责明视觉、颜色与较精细的空间细节；视杆在弱光下更敏感。人眼不同位置的感受器分布不同，不能把整个视网膜看作性能完全一致的相机阵列。
\item \textbf{亮度适应：}人眼可适应很宽的光强范围，但不能在同一适应状态下同时分辨整个范围。进入暗室后逐渐看清，不是物体自身突然变亮。
\item \textbf{亮度辨别：}在背景强度 $I$ 上增加 $\Delta I$，能否看出变化取决于背景与视觉状态。韦伯比为
\[
W=\frac{\Delta I_{50}}{I},
\]
其中 $\Delta I_{50}$ 表示实验中有50\%概率被辨别出的增量。\textbf{比值越小，说明该条件下对相对变化越敏感。}
\item \textbf{马赫带：}实际强度分段变化时，人会在边界附近看到额外的亮暗带；感知亮度不是像素强度的直接复制。
\item \textbf{同时对比：}相同灰度的小块，放在暗背景上往往显得更亮，放在亮背景上往往显得更暗。评价增强结果时要考虑邻域背景。
\end{itemize}

\paragraph{自编例题S1：韦伯比与“能否分辨”}
实验A的背景 $I_A=100$，阈值增量 $\Delta I_{50,A}=2$；实验B的背景 $I_B=1000$，阈值增量为30。哪个条件下的相对亮度辨别能力更好？若在同一适用范围内假设 $W=0.02$，背景200需要多大阈值增量？

\textbf{解：}
\[
W_A=2/100=0.02,\qquad W_B=30/1000=0.03.
\]
A的相对辨别能力更好。按题设近似，新的阈值增量为 $0.02\times200=4$。
\textbf{注意：}不能仅比较2和30两个绝对增量；也不能把韦伯比在所有光强下都当作固定常数。

\paragraph{易错判断}
“图像中两个区域灰度相等，所以人眼一定觉得它们同样亮。”\textbf{错误。}背景、适应状态等会影响主观亮度。“灰度级越多，所有相邻灰度都一定能看出差别。”也错误，显示条件与辨别阈值同样重要。

\clearpage
\subsection{像素邻接、连通性与距离}
\textit{来源：Sampling 课件第26--28页；教材\S2.5，印刷页79--83（PDF页81--85）。距离公式为教材补充。}

对像素 $p=(x,y)$：
\[
\begin{aligned}
N_4(p)&=\{(x-1,y),(x+1,y),(x,y-1),(x,y+1)\},\\
N_D(p)&=\{(x-1,y-1),(x-1,y+1),(x+1,y-1),(x+1,y+1)\},\\
N_8(p)&=N_4(p)\cup N_D(p).
\end{aligned}
\]
图像边界外的位置不算有效像素。邻接还要指定允许的灰度集合 $V$，例如二值前景取 $V=\{1\}$。只有满足灰度条件的像素，才参与相应前景连通关系。
\begin{itemize}
\item \textbf{4邻接：}$q\in N_4(p)$，并且两像素灰度均在 $V$ 中。
\item \textbf{8邻接：}$q\in N_8(p)$，并且两像素灰度均在 $V$ 中。
\item \textbf{混合邻接（m邻接）：}两像素灰度均在 $V$ 中，且或者4邻接；或者为对角邻居，同时 $N_4(p)\cap N_4(q)$ 中\textbf{没有灰度属于 $V$ 的像素}。这个附加条件用于避免部分对角连接的歧义。
\item \textbf{连通：}能通过指定邻接规则的路径到达；一个连通分量是按该规则相互可达的最大像素集合。邻接是一步关系，连通允许走多步。
\end{itemize}
对 $p=(x,y)$、$q=(u,v)$，常见坐标距离为
\[
D_E=\sqrt{(x-u)^2+(y-v)^2},\quad
D_4=|x-u|+|y-v|,\quad
D_8=\max(|x-u|,|y-v|).
\]
没有障碍的网格上，$D_4,D_8$ 分别对应最短4邻接、8邻接步数；有前景限制时，实际可行路径可能更长，甚至不存在。

\paragraph{自编例题S2：对角像素一定m邻接吗？}
设二值图像为
\[
\begin{bmatrix}1&1\\0&1\end{bmatrix},\qquad
p=(0,0),\quad q=(1,1),\quad V=\{1\}.
\]
问 $p,q$ 是否4邻接、8邻接、m邻接？是否4连通？

\textbf{解：}两者是对角邻居，因此不是4邻接，但为8邻接。共同4邻域包含 $(0,1)$ 与 $(1,0)$；其中 $(0,1)$ 的值为1，所以\textbf{不是m邻接}。然而有路径
\[
(0,0)\longrightarrow(0,1)\longrightarrow(1,1),
\]
所以它们仍然4连通，也m连通。相应坐标距离为 $D_E=\sqrt2$、$D_4=2$、$D_8=1$；本图的最短m路径有2步。\textbf{若把右上角的1改为0，}它们变成m邻接，但不再4连通。

\clearpage
\subsection{多帧平均去噪：为何有效，何时失效}
\textit{来源：增强课件第29、34、54--57页；教材\S2.6，印刷页85--88（PDF页87--90）。}

将同一场景的第 $i$ 帧写成 $g_i=f+n_i$，平均结果为
\[
\bar g=\frac1N\sum_{i=1}^{N}g_i=f+\frac1N\sum_{i=1}^{N}n_i.
\]
若噪声零均值、各帧之间不相关且每帧噪声方差均为 $\sigma^2$，则
\[
E[\bar g]=f,\qquad
\operatorname{Var}(\bar n)=\frac{\sigma^2}{N},\qquad
\sigma_{\bar n}=\frac{\sigma}{\sqrt N}.
\]
\textbf{需要同时检查：}场景不变或已经配准；亮度一致；噪声近似零均值且帧间不相关。物体运动会导致重影，固定图案噪声和系统偏差不会因简单平均自动消失。

\paragraph{自编例题S3：需要平均多少帧？}
单帧噪声标准差为20灰度级，欲将平均图噪声标准差降到不超过4，至少要多少帧？

\textbf{解：}
\[
\frac{20}{\sqrt N}\le4
\ \Longrightarrow\ \sqrt N\ge5
\ \Longrightarrow\ N\ge25.
\]
25帧时方差由400降到16，标准差由20降到4。\textbf{不是}标准差除以25。若只平均4帧，标准差为10。

\paragraph{进一步理解：相关噪声为何难消除？}
一般形式是
\[
\operatorname{Var}(\bar n)=\frac1{N^2}
\left[\sum_i\operatorname{Var}(n_i)+2\sum_{i<j}\operatorname{Cov}(n_i,n_j)\right].
\]
只有协方差项为零时，才直接得到 $\sigma^2/N$。若各帧均有同样的加性偏差 $b$，平均后仍保留 $b$。本式用于理解假设，不作为已确认的推导考点。

\paragraph{与空间均值滤波的区别}
多帧平均在\textbf{同一空间位置跨时间}平均，静态且对齐时可降噪而不必模糊空间边缘。空间均值滤波在\textbf{单帧邻域内}平均，通常会同时模糊边缘。两者不能仅因名字都有“平均”就混为一谈。

\paragraph{实现提醒}
累加器应初始化为零，累加 $N$ 帧后除以 $N$；中间运算使用足够精度，最后再裁剪和转换数据类型。增强课件第29页若先把累加器设为原图，再累加50帧并除以50，会多算一份原图，应修正。该提醒是代码校对，不是额外押题。


\subsection{课件扩展：BBHE与亮度保持}
增强课件第59--65页：BBHE按输入均值拆分直方图，分别均衡，旨在缓解亮度漂移，但不保证精确保持均值。
用 $\mathrm{AMBE}=|\mu_{\rm out}-\mu_{\rm in}|$ 衡量均值变化。例如输入均值80，输出125或86，则AMBE为45或6；后者亮度更接近，却不代表细节与噪声一定更好。
课件中的精确保持条件依赖理想连续模型，不能直接当作一般离散图像恒等式。
